The Hardware Layer comprises \textbf{bench capture endpoints}: Raspberry Pi devices running
Agri-Home/klipper with \texttt{fswebcam} camera macros and the AgriHome edge agent. It is the
physical source of tray imagery for automatic and operator-triggered capture. Unlike the
Presentation Layer, these nodes authenticate with hashed \textbf{device API keys}, not Firebase
user sessions.

\subsection{Raspberry Pi edge host subsystem}

Each provisioned \textbf{Raspberry Pi} (or compatible board) is a durable edge host with a CPU serial fingerprint used at registration. AgriHome stores the device in \texttt{edge\_devices}, links it to a tray, and issues a one-time API key that the agent persists locally. Register uses \texttt{POST /api/raspberry-pi/register} with \texttt{DEVICE\_PROVISIONING\_SECRET}; the default model string is \texttt{raspberry-pi-klipper}.

\begin{figure}[H]
\centering
\begin{tikzpicture}[font=\scriptsize, box/.style={draw, rounded corners, align=center, minimum width=2.3cm, minimum height=0.7cm}, arr/.style={-{Stealth}, thick}]
\node[box, fill=teal!14] (pi) {Raspberry Pi\\host};
\node[box, right=1.0cm of pi, fill=blue!8] (api) {Application\\APIs};
\draw[arr] (pi) -- node[above] {register / heartbeat} (api);
\end{tikzpicture}
\caption{Pi edge host: register and heartbeat against AgriHome.}
\label{fig:hw-pi}
\end{figure}

\subsubsection{Assumptions}
\begin{itemize}
  \item The host has network reachability to the AgriHome base URL.
  \item A shared provisioning secret is configured on both server and agent.
\end{itemize}

\subsubsection{Responsibilities}
\begin{itemize}
  \item Present stable hardware identity (CPU serial, MAC, hostname, model).
  \item Store the device API key securely on disk after first registration.
  \item Remain online via periodic heartbeats so Vision Console status stays current.
\end{itemize}

\subsubsection{Subsystem interfaces}
\begin{table}[H]
\centering
\small
\begin{tabular}{|p{0.9cm}|p{3.4cm}|p{3.0cm}|p{3.5cm}|}
\hline
\textbf{ID} & \textbf{Description} & \textbf{Inputs} & \textbf{Outputs} \\
\hline
H-P1 & Pi host $\rightarrow$ Application & Fingerprint + provisioning code & \texttt{deviceId}, \texttt{trayId}, API key \\
\hline
H-P2 & Pi host $\rightarrow$ Application & Device-key heartbeat & Online status; pending commands \\
\hline
\end{tabular}
\caption{Raspberry Pi edge host subsystem interfaces}
\end{table}

\subsection{Klipper / fswebcam capture subsystem}

\textbf{Agri-Home/klipper}\cite{klipper} provides G-code/macros for axis motion and
\texttt{camera-macros/save\_image.sh}, which captures JPEG stills with \texttt{fswebcam}. This is
the \textbf{primary} Take Picture path when the agent claims \texttt{capture\_now}. An optional LAN
HTTP streamer base may be stored as \texttt{klipper\_url} (migration \texttt{012\_klipper\_url} renamed
\texttt{moonraker\_url}$\rightarrow$\texttt{klipper\_url}) so the Application Layer can attempt a
server-side still when reachable. Register/heartbeat accept deprecated \texttt{moonrakerUrl} as an
alias for \texttt{klipperUrl}.

\begin{figure}[H]
\centering
\begin{tikzpicture}[font=\scriptsize, box/.style={draw, rounded corners, align=center, minimum width=2.2cm, minimum height=0.7cm}, arr/.style={-{Stealth}, thick}]
\node[box, fill=blue!8] (app) {Application};
\node[box, right=1.0cm of app, fill=teal!12] (cam) {fswebcam\\save\_image.sh};
\node[box, below=0.85cm of cam, fill=teal!10] (agent) {agrihome\_agent};
\node[box, left=1.0cm of agent, fill=gray!12, dashed] (str) {optional\\HTTP streamer};
\draw[arr] (agent) -- (cam);
\draw[arr, dashed] (app) -- node[above] {GET still} (str);
\end{tikzpicture}
\caption{Primary fswebcam capture on the agent; optional server-side HTTP streamer still.}
\label{fig:hw-moonraker}
\end{figure}

\subsubsection{Assumptions}
\begin{itemize}
  \item \texttt{fswebcam} and Agri-Home/klipper camera macros are installed on the Pi.
  \item Optional server-side Take Picture requires the AgriHome host to reach the streamer URL on the LAN (WARP local-network overrides may be required).
\end{itemize}

\subsubsection{Responsibilities}
\begin{itemize}
  \item Produce JPEG stills for operator and scheduled capture via \texttt{save\_image.sh}.
  \item Optionally expose an HTTP streamer base for operators who want an immediate server pull.
  \item Keep \texttt{klipper\_url} as an optional streamer/API base distinct from the agent capture command.
\end{itemize}

\subsubsection{Subsystem interfaces}
\begin{table}[H]
\centering
\small
\begin{tabular}{|p{0.9cm}|p{3.4cm}|p{3.0cm}|p{3.5cm}|}
\hline
\textbf{ID} & \textbf{Description} & \textbf{Inputs} & \textbf{Outputs} \\
\hline
H-M1 & Application $\rightarrow$ streamer & Snapshot URL candidates & JPEG bytes or timeout \\
\hline
H-M2 & agrihome\_agent $\rightarrow$ fswebcam & \texttt{save\_image.sh} path & JPEG bytes \\
\hline
\end{tabular}
\caption{Klipper / fswebcam capture subsystem interfaces}
\end{table}

\subsection{agrihome\_agent bridge subsystem}

The \textbf{\texttt{agrihome\_agent}} Python package (still hosted under the Agri-Home/moonraker companion repo, branch \texttt{feature/agrihome-bridge}) registers the device, heartbeats, claims \texttt{capture\_now} commands, walks pose sequences when requested, runs \texttt{fswebcam} via Klipper camera macros, and multipart-ingests frames to AgriHome. It speaks \texttt{klipperUrl} (with deprecated \texttt{moonrakerUrl} alias).

\begin{figure}[H]
\centering
\begin{tikzpicture}[font=\scriptsize, box/.style={draw, rounded corners, align=center, minimum width=2.3cm, minimum height=0.7cm}, arr/.style={-{Stealth}, thick}]
\node[box, fill=teal!12] (agent) {agrihome\_agent};
\node[box, right=1.0cm of agent, fill=blue!8] (api) {/api/raspberry-pi/*};
\node[box, below=0.85cm of agent, fill=teal!10] (cam) {fswebcam / Klipper};
\draw[arr] (agent) -- node[above] {ingest / claim} (api);
\draw[arr] (agent) -- (cam);
\end{tikzpicture}
\caption{Edge agent bridges Klipper/fswebcam local capture to AgriHome APIs.}
\label{fig:hw-agent}
\end{figure}

\subsubsection{Assumptions}
\begin{itemize}
  \item Agent configuration lives under a local agrihome agent config file (or equivalent environment variables), including \texttt{AGRIHOME\_SNAPSHOT\_CMD} pointing at \texttt{save\_image.sh}.
  \item Failed individual poses do not abort the remainder of a multi-pose run.
\end{itemize}

\subsubsection{Responsibilities}
\begin{itemize}
  \item Execute register / run loops against documented \texttt{/api/raspberry-pi/*} routes.
  \item Claim and complete edge commands; upload multipart captures with pose metadata.
  \item Fetch active pose sequences and capture plans for scheduled multi-angle work.
\end{itemize}

\subsubsection{Subsystem interfaces}
\begin{table}[H]
\centering
\small
\begin{tabular}{|p{0.9cm}|p{3.4cm}|p{3.0cm}|p{3.5cm}|}
\hline
\textbf{ID} & \textbf{Description} & \textbf{Inputs} & \textbf{Outputs} \\
\hline
H-A1 & Agent $\rightarrow$ Application & Device-key HTTP & Commands, plans, ingest ack \\
\hline
H-A2 & Agent $\rightarrow$ Klipper/fswebcam & Snapshot / motion cues & Frame bytes; pose dwell \\
\hline
\end{tabular}
\caption{\texttt{agrihome\_agent} bridge subsystem interfaces}
\end{table}
